% ============================
%           GLOSARIO
% ============================

\chapter*{Glosario}

\addcontentsline{toc}{chapter}{Glosario}

\begingroup
\small

\renewcommand{\arraystretch}{1.15}

\begin{longtable}[l]{
    @{}
    >{\raggedright\arraybackslash}p{4.2cm}
    p{\dimexpr\textwidth-4.2cm\relax}
    @{}
    }

    \toprule
    \textbf{Término} & \textbf{Definición}                                \\
    \midrule
    \endfirsthead

    \toprule
    \textbf{Término} & \textbf{Definición}                                \\
    \midrule
    \endhead

    \midrule
    \multicolumn{2}{r}{\small\textit{Continúa en la siguiente página...}} \\
    \endfoot

    \bottomrule
    \endlastfoot

    Acción legal
                     &
    Movimiento permitido para un jugador en un estado específico del juego,
    de acuerdo con las reglas declaradas en la descripción formal del dominio.
    \\

    Agente inteligente
                     &
    Sistema computacional capaz de percibir un entorno, razonar sobre sus
    posibles acciones y seleccionar una decisión con base en un objetivo,
    recompensa o criterio de utilidad.
    \\

    Árbol de juego
                     &
    Representación de las posibles secuencias de estados y acciones de un
    juego. Cada nodo representa un estado y cada arista representa una acción
    que conduce a un nuevo estado.
    \\

    Descripción de juego
                     &
    Especificación externa que define jugadores, estado inicial, acciones
    legales, transiciones, estados terminales y recompensas sin codificar
    directamente esas reglas dentro del motor del agente.
    \\

    Estado
                     &
    Configuración completa de un juego en un instante determinado. Debe
    contener la información necesaria para determinar acciones legales,
    transiciones y condiciones terminales.
    \\

    Estado terminal
                     &
    Estado en el cual la partida termina y se asigna una recompensa o utilidad
    a cada jugador.
    \\

    Función heurística
                     &
    Función de evaluación que estima qué tan favorable es un estado cuando no
    resulta viable explorar el árbol de juego hasta estados terminales.
    \\

    General Game Player
                     &
    Agente capaz de recibir una descripción formal de un juego y jugarlo sin
    que su motor de decisión contenga reglas específicas programadas para ese
    dominio.
    \\

    Información perfecta
                     &
    Propiedad de los juegos en los que todos los jugadores conocen el estado
    completo del juego y no existen reglas ocultas para alguno de ellos.
    \\

    Intérprete
                     &
    Componente encargado de analizar el archivo de reglas, validar su sintaxis
    y semántica, y construir una representación interna que pueda consultar el
    motor del agente.
    \\

    Lenguaje de descripción de juegos
                     &
    Lenguaje formal o de dominio específico utilizado para expresar reglas de
    juegos de manera declarativa y separada del código del motor.
    \\

    Motor de decisión
                     &
    Módulo encargado de explorar estados posibles y seleccionar una acción
    mediante algoritmos generales de búsqueda o simulación.
    \\

    Poda alfa-beta
                     &
    Optimización del algoritmo Minimax que evita explorar ramas del árbol de
    juego que no pueden cambiar la decisión final.
    \\

    Recompensa
                     &
    Valor asignado a un jugador, generalmente al alcanzar un estado terminal,
    que indica el resultado o utilidad de la partida para ese jugador.
    \\

    Transición
                     &
    Cambio de un estado a otro como resultado de aplicar una acción legal.
    \\

\end{longtable}

\endgroup
